% ! TEX root = ../../main.tex

% Prev intro: Checkpoint post-processing is bad In-situ exists but no range queries Our goal:
% reordering-based in-situ indexing mechanisms CARP and how it solves the problem

CARP enables ingestion of scientific checkpoint data for efficient range queries in a single
streaming pass, via an adaptive in-situ range partitioning mechanism.
Range queries over continuous attributes like temperature and velocity are fundamental to
scientific analysis, enabling filtering of high-energy particles~\cite{Byna2012} or tracking
shock waves in fluid simulations~\cite{Grete2023:Parthenon, Zhang2021:AMReX, Barke2024:Phoebus}.
Efficient range queries require key locality in the stored layout, but existing mechanisms
for creating such layouts require post-processing.
This incurs additional I/O passes and reduces
effective ingestion bandwidth.
As checkpoint data volumes grow with application scale, these post-processing costs increasingly
dominate time-to-discovery~\cite{Bauer2016:SituMethodsInfrastructures}.

Among existing indexing approaches, those that enable range queries force undesirable tradeoffs
between
ingestion and query performance, while in-situ approaches do not enable range queries.
Distributed out-of-core sorts produce fully ordered layouts which enable both efficient
\emph{location}
of overlapping keys and efficient \emph{retrieval}, via large sequential reads.
This enables excellent
query performance, but requires two additional read and write passes over the data.
Auxiliary bitmap indices do not reorder data but generate additional index structures.
They can accelerate location of matching keys, but retrieval still
incurs additional random reads over scattered keys.
Database-style ingestion is unviable: it is designed for mixed workloads but
write amplification from periodic reorganization makes it impractical for ingestion-heavy
scientific workloads~\cite{Raju2017:PebblesDB}.
Among in-situ approaches, DeltaFS~\cite{Zheng2018:DeltaFS} demonstrates the viability of in-situ
hash partitioning
via an embedded all-to-all shuffle on MPI ranks, but hash-based
partitioning destroys key locality and does not support range queries.

The central challenge of extending in-situ partitioning to range queries is that of providing
efficient ingestion for unknown key distributions.
Our analysis of real scientific codes shows that these distributions are often highly skewed and
can shift significantly across epochs.
Distributed databases manage range-partitioned layouts by reactively reordering on-disk data
in response to partition imbalance, incurring unacceptable write amplification.
Our goal is an adaptive partitioner that makes no assumptions about key distributions,
and optimizes for ingestion performance by never reordering on-disk data, while still enabling
efficient range queries.
This requires mechanisms for discovering global key distributions,
reacting as they evolve, handling out-of-bounds keys, all while maintaining high
storage utilization.

% Boundaries computed from early data become stale as distributions evolve, creating load hotspots
% during ingestion and degraded selectivity during retrieval. Adapting to these shifts during a live
% checkpoint introduces a design problem: the system must decide when to repartition, how to
% construct a current global view of the distribution from per-rank local state, how to compute new
% boundaries from that view, and how to handle keys that fall outside existing partition ranges as
% distributions drift. These concerns interact---repartitioning frequency affects view staleness,
% view quality affects boundary balance, and boundary changes affect how in-flight data is
% routed---and the entire mechanism must operate at negligible cost relative to checkpoint I/O.

CARP realizes adaptive single-pass range partitioning via a feedback loop for discovering
and reacting to global key distributions.
A primitive called \emph{renegotiation} is used to materialize the view for this loop: a reduction
tree
over local histograms maintained at shuffle senders materializes the current global distribution,
which is used to compute and broadcast new load-balanced partition boundaries.
% Called \emph{renegotiation}, it uses a reduction tree over local histograms maintained at shuffle
% senders to materialize the current global distribution, which is then used to compute and
% broadcast new load-balanced partition boundaries.
Trigger mechanisms invoke this process periodically, and these partition boundaries inform
shuffling until the next renegotiation.
An out-of-bounds mechanism enables both bootstrapping and handling keys that fall outside the
current partition ranges.
Shuffle receivers accumulate writes to per-rank append-only logs, never reordering on-disk data.
Because the storage layer is fully utilized with append-only I/O,
CARP incurs no write overhead over unpartitioned I/O in evaluated configurations.
The resulting data layout is partially ordered, which requires query-time reconciliation.
In our evaluation, this reconciliation overhead is negligible,
with query performance comparable to fully sorted layouts and upto $4.9\times$ faster
ingestion.

The rest of this chapter provides background on data organization for range queries and the
tradeoffs involved (\S\ref{sec:carp:bg}), characterizes key distributions and their drift in real
scientific codes (\S\ref{sec:carp:distrib}), presents the design of CARP and its storage backend
KoiDB (\S\ref{sec:carp:des}), evaluates write and query performance (\S\ref{sec:carp:eval_write},
\S\ref{sec:carp:eval_read}), and discusses adaptations for different deployment requirements
(\S\ref{sec:carp:disc}).
%The CARP framework and KoiDB backend are open source~\cite{Jain2024:CARP}.
